\scanpage{15}
\mspage{13}
12) Structures supplémentaires sur le groupe de Galois motivique : \inserted{structures associées aux} graduations, filtrations, au motif de Tate.

\deleted{La catégorie} On choisit une sous-\inserted{$\otimes$-}catégorie graduée $C$ de $M(S)$, avec un foncteur fibre $F$, d'où groupe de Galois motivique $G$

a) Une graduation de type $\unsure{\pi|\Pi|M}$ sur le foncteur fibre \inserted{à coeff dans $K$} compatible avec $\otimes$, est équivalent à la donnée d'\unsure{une|un} hom
\[
T=D_K(\pi)\longrightarrow G
\]
La graduation provient d'une graduation fonctorielle de $M\in\inserted{\operatorname{Ob}C}$ lui-\unsure{$=$|même} \unsure{sss|ssi} l'hom $T\longrightarrow G$ est central. Par suite :

$\underline{1^\circ})$ La graduation \inserted{canonique} de type $\mathrm Z$ des objets de $C$ définit un hom canonique
\[
\unsure{\eta|\gamma|y}:\mathrm{G}_{m\,K}\longrightarrow G
\]
$2^\circ)$ Si $F=T_{\mathrm{Hdg}}$, la bigraduation de ce foncteur correspond à un hom
\[
\deleted{D_K(\mathrm Z^2)=\mathrm G_m^2\longrightarrow G,\quad\text{donc à deux h}}
\]
\[
i:\mathrm G_{m\,K}^2\longrightarrow G
\]
donc à deux hom commutant l'un à l'autre
\[
i_1,i_2:\mathrm G_m\unsure{\longrightarrow|-}G
\]
Comme la graduation totale associée est celle de $1^\circ)$, on trouve
\[
i_1(\lambda)i_2(\lambda)=\eta(\lambda)
\]
(ce qui entraîne déjà que $i_1,i_2$ commutent\ldots).
% archive sheet number: 14

